\subsection*{A genuine gentle action on the quotient by the rank-129 kernel}

\paragraph{GD.1 (Proved): the finite two-group inputs and their central ambiguity.}
Retain $D=C''(C')^2$, $Q=C/D$ and $\Gamma_D=L/D$.
In the exact cube-chain model of $Q$, the images in
$\operatorname{Aut}(Q)$ of the five pairs have nilpotency
class at most two. The two harmonic pair groups have
order eight and central derived subgroup of order two;
the other three pairs commute. In particular every
$f\in\gamma_3(F)$ evaluates to one in these five
automorphism-pair groups. This is checked directly
from the literal involution-word actions, independently
in \texttt{check\_gentle\_Q\_action.py}: each pair
commutator commutes with both generators. The input
is the actual $256$-element quotient, not a parameter
assignment for $f$.

The following central facts can also be seen directly
in the cube. Write $S_i$ for the sum, over $\mathbb F_2$,
of the four edges in direction $i$. The cycle space
$C'/D$ has dimension five. The action of
$C/C'=(C_2)^3$ on it is faithful: a nonzero translation
in direction $i$ exchanges the two disjoint face
cycles normal to that direction. Therefore a central
element of $Q$ has trivial image in $C/C'$. The
translation-invariant edge chains are the span of
$S_0,S_1,S_2$, and each $S_i$ has boundary the sum
of all eight vertices. Thus
\[
 Z(Q)=\{c_0S_0+c_1S_1+c_2S_2:c_0+c_1+c_2=0\}
           \simeq(C_2)^2\subset Q'.
 \tag{GD.1a}
\]
The exact path substitutions for $a$ interchange
$S_0,S_1$ and fix $S_2$; those for $b$ interchange
$S_0,S_2$ and fix $S_1$. A common invariant vector
in the displayed even-sum plane is therefore zero.
An element of $Q$ fixed by $L$ first centralizes
$C\subset L$, so belongs to this plane. Consequently
\[
                              Q^L=1.
 \tag{GD.1b}
\]
Both assertions are also verified in the independent
finite certificate. In its twelve-edge binary
encoding the three nonzero centre vectors are
$255,3855,4080$.

\paragraph{GD.2 (Proved): all relators, kernel stability, and the composition law on $L/D$.}
Let $\Phi=(1,f)$ be a genuine gentle element. Its
two- and three-cycle relations imply
$f\in\gamma_3(F)$, as proved in RP.8. Write
\[
 g=f(X_1,X_2),\quad j=f(X_2,X_3),\quad
 h=f(X_4,X_3),\quad k=f(X_5,X_1),\quad l=f(X_4,X_5).
\]
All words in this paragraph are evaluated in the
actual quotient $\Gamma_D$. Define generator images
\[
              a'=a,\qquad b'=g^{-1}bg,
                              \qquad c'=h^{-1}ch.
 \tag{GD.2a}
\]
The first braid relation for $a',b'$ is the genuine
three-strand gentle relation; its proof uses only
the two- and three-cycle equations and the source
free-group automorphism. It is the three-strand
formula recorded in FP.26.4, in the same marking.
For clarity about the central lift, the gentle
extension on the central quotient fixes the
projective braid relation; its residual central
error upstairs has exponent sum zero and hence
is one. The image of $z_3=(ab)^3$ centralizes the
image of the free pure kernel, which is still
that kernel by the genuine free-group
automorphism. It therefore belongs to the
three-strand centre. If it is $z_3^t$, exponent
sum gives $6t=6$ in $\widehat{\mathbb Z}$,
so $t=1$. Thus $z_3$ is fixed, without a
five-point action assumption.

The commutation $[a',c']=1$ holds because $a$
commutes with both arguments $X_4=z_3$ and
$X_3=c^2$ of $h$. For the remaining braid relation,
RP.8 gives the actual quotient identity
\[
                   h^{-1}kjlg=1,
                   \qquad h=kjlg.
 \tag{GD.2b}
\]
Here $k$ commutes with $c$, since both $X_5=(bc)^3$
and $a^2$ commute with $c$. Likewise $l$ commutes
with $b$, since both three-strand centres commute
with $b$. Conjugating $b',c'$ by $g$ therefore
gives
\[
       b,\qquad (hg^{-1})^{-1}c(hg^{-1})
                  =l^{-1}j^{-1}cjl.
\]
The genuine three-strand gentle relation for
$b,j^{-1}cj$, followed by conjugation by $l$,
proves the desired braid relation. Thus the
assignment already defines a homomorphism
$\widehat B_4\to\Gamma_D$.

Its image of the central full twist $z_4=(abc)^4$
has trivial fold in $\widehat B$, because the
folded generator images are the genuine gentle
three-strand images of $a,b,a$. Hence this full
twist image belongs to $Q$. By GD.1, $g$ and $h$
centralize $Q$, so the new generators have exactly
the old conjugation actions on $Q$. The full
twist image centralizes all new generators,
and therefore belongs to $Q^L=1$. This verifies
the fourth defining relator, yielding a genuine
continuous homomorphism
\[
                    \widetilde\Phi:L\longrightarrow\Gamma_D.
\]

Its fold is the original gentle automorphism of
$\widehat B$. Thus $\widetilde\Phi(C)\subset Q$.
For $w\in C$, its new and old images induce the
same inner action on $Q$, because the generator
actions agree. Consequently
\[
       \delta_\Phi(w)=\widetilde\Phi(w)\bar w^{-1}\in Z(Q).
 \tag{GD.2c}
\]
This is a homomorphism $C\to Z(Q)$: the differences
are central, so the multiplicative error for a
product is the product of the errors. It kills
$C'$, and therefore kills $D\subset C'$.
Since $\bar D=1$, the map $\widetilde\Phi$
kills $D$. This proves the required kernel
stability before defining
\[
                    \Phi_D:\Gamma_D\longrightarrow\Gamma_D.
\]

The maps obey the genuine gentle composition law.
Use the convention $\Phi_{12}=\Phi_1\circ\Phi_2$,
so on the free source
$f_{12}=f_1\Phi_1(f_2)$. On the first pair
$(a^2,b^2)$, the map $\Phi_{1,D}$ fixes the first
argument and conjugates the second by
$g_1=f_1(a^2,b^2)$. Thus
$g_{12}=g_1\Phi_{1,D}(g_2)$. The same calculation
on $(z_3,c^2)$, using $\Phi_{1,D}(z_3)=z_3$,
gives $h_{12}=h_1\Phi_{1,D}(h_2)$.
Substitution in (GD.2a) proves
\[
              \Phi_{12,D}=\Phi_{1,D}\circ\Phi_{2,D}.
\]
The identity pair gives the identity map, and
the inverse gentle pair supplies the inverse.
We have therefore proved an actual continuous
action
\[
                 \boxed{\ \mathcal H\longrightarrow
                             \operatorname{Aut}(L/D).\ }
 \tag{GD.2d}
\]
All defining relators and the kernel have been
checked; this is not an action inferred solely
from the vanishing of a pentagon product.

\paragraph{GD.3 (Proved): the entire gentle group fixes $Q$ pointwise, with no inferred action on $M$.}
The correction (GD.2c) factors through
$C^{\rm ab}=(\mathbb F_2)^3$ and is $L$-equivariant,
because the new and old $L$ conjugation actions
on $Q$ agree. The action on $Q'$ is pointwise
the identity. Since $Z(Q)\subset Q'$, composition
adds these corrections. Their possible values form
\[
       \operatorname{Hom}_L((\mathbb F_2)^3,Z(Q))
                              \simeq\mathbb F_2.
 \tag{GD.3a}
\]
Indeed the first module is the three-letter
permutation module, and $Z(Q)$ is its even-sum
two-dimensional summand with the same $S_3$ action.
The one-dimensional Hom-space is also checked
by enumerating its sixty-four possible linear
maps in the exact certificate. This is only
an intermediate bound: the actual correction
always vanishes, as follows.

In $\Gamma_D$, let $J=\overline{\langle z_3,c^2\rangle}$.
Its fold is procyclic, so $J'\subset Q$.
Its conjugation action on $Q$ is abelian
by GD.1, so in fact $J'\subset Z(Q)$.
Both generators act trivially on $Z(Q)$:
$z_3\in Q$, and $c^2$ has trivial
three-letter permutation action on the
even-sum plane (GD.1a). Thus $J$ is
nilpotent of class at most two. Since
$f\in\gamma_3(F)$, this proves
\[
                        h=f(z_3,c^2)=1
                         \quad\text{in }\Gamma_D.
\]
The same cyclic-fold and central-action
argument applies to the other two
nonharmonic pairs, giving $k=l=1$.
In particular (GD.2a) fixes $a$ and $c$,
and therefore fixes $u=ca^{-1}=t_1t_0$.
It follows that
$\delta_\Phi(t_0)=\delta_\Phi(t_1)=z$
for some $z\in Z(Q)$. Equivariance under
$a$, which interchanges $t_0,t_1$ in
the abelianization, gives $a(z)=z$.
Equivariance under $b$, which fixes
$t_1$, gives $b(z)=z$. Equation (GD.1b)
forces $z=0$. Finally $b$ interchanges
$t_0,t_2$ in the abelianization, so
$\delta_\Phi(t_2)=0$ as well. Therefore
\[
          \boxed{\ \Phi_D|_{Q}=\operatorname{id}_Q
                     \quad\text{for every }\Phi\in\mathcal H.\ }
 \tag{GD.3b}
\]
There is no nontrivial central-correction
character produced by this quotient.
The restriction is trivial on the entire
gentle group, including the FR.7 witness
with $\Delta=1$.

An action on $L/D$ does not, however, determine
an action on $M=H_1(D,\mathbb F_3)$, because
$D$ has already been killed. A lift through
the free-product cover $C\to Q$ requires
peripheral conjugators, and changing those
lifts can change their action on $M$.
The argument above does not construct such
a lift for all gentle elements, nor show
that $k_{\rm rel}$ acts trivially on $M$.
Consequently it supplies no crossed-cocycle
composition law for
$m_f=[\mathcal P(f)]\in D/(D'D^3)$.
Such a law requires an additional specified
action on this module and a compatibility
proof. In particular no three-power index
or odd-power witness argument is inferred
from an unconstructed affine action.
